\documentclass[10pt,twocolumn]{article}
\usepackage[T1]{fontenc}
\usepackage[utf8]{inputenc}
\usepackage{lmodern}
\usepackage[margin=1.8cm,top=2.4cm,bottom=2cm,columnsep=0.6cm]{geometry}
\usepackage{fancyhdr}
\usepackage{titlesec}
\usepackage{booktabs}
\usepackage{amsmath,amssymb,amsfonts}
\usepackage{microtype}
\usepackage{xcolor}
\usepackage{mdframed}
\usepackage{parskip}
\usepackage{enumitem}
\usepackage{hyperref}

\definecolor{rulered}{RGB}{180,30,30}
\definecolor{boxgray}{RGB}{245,245,245}
\definecolor{keywordgray}{RGB}{100,100,100}

\pagestyle{fancy}
\fancyhf{}
\fancyhead[L]{\textsc{\small Claude Mag}}
\fancyhead[C]{\small\textit{Machine Learning Research Digest}}
\fancyhead[R]{\small 2026-09-09}
\fancyfoot[C]{\small\thepage}
\renewcommand{\headrulewidth}{0.8pt}

\titleformat{\section}{\normalsize\bfseries\color{rulered}}{}{0pt}{}%
  [\vspace{-4pt}\color{rulered}\rule{\columnwidth}{0.4pt}]
\titlespacing{\section}{0pt}{8pt}{4pt}

\hypersetup{colorlinks=true,urlcolor=rulered,linkcolor=black}

\begin{document}
\setlength{\parindent}{0pt}
\setlength{\parskip}{4pt}

{\small\color{keywordgray}\textbf{Keywords:}
  multimodal LLM \textbullet{}
  hallucination \textbullet{}
  visual grounding \textbullet{}
  object hallucination}\par\vspace{4pt}

{\LARGE\bfseries\raggedright
  Beyond Language Priors: Diagnosing and Fixing
  Visual-Origin Hallucinations in Multimodal LLMs}\par\vspace{4pt}

{\large\itshape\raggedright
  Cosine Similarity and Grad-CAM Entropy Reveal
  That Half of Object Hallucinations Originate
  in the Visual Encoder, Not the Language Model}\par\vspace{6pt}

{\small\textbf{By} Peiyang Xu, Xiaopei Zhu, Jun Zhu, Xiaolin Hu
  (Tsinghua University)}\par\vspace{2pt}

{\small\color{keywordgray}arXiv:2609.00231 \textbullet{} Preprint, August 2026}
\par\vspace{2pt}\noindent{\color{rulered}\rule{\columnwidth}{0.6pt}}\vspace{6pt}

Object hallucination in multimodal large language models (MLLMs) --- the
generation of text describing objects not present in the image --- is
a persistent failure mode that undermines trust in deployed systems.
The field has converged on a near-consensus explanation: the language
model component over-relies on learned co-occurrence statistics (language
priors), producing objects that are likely in the described scene even
when they are absent from the image. This paper challenges that consensus.
Using cosine similarity analysis and Smooth Grad-CAM entropy, the authors
show that 38--52\% of object hallucinations across three MLLM architectures
originate not in the language prior but in failures of the visual encoder
itself.

\begin{mdframed}[backgroundcolor=boxgray,linecolor=rulered,linewidth=1pt,
    innerleftmargin=6pt,innerrightmargin=6pt,
    innertopmargin=6pt,innerbottommargin=6pt]
\textbf{\small AT A GLANCE}\par\vspace{4pt}
\begin{description}[leftmargin=0pt,labelsep=4pt,itemsep=2pt,parsep=0pt,
    font=\small\bfseries]
  \item[Problem:] Object hallucination in MLLMs is blamed on language
    priors; visual encoding failures are ignored.
  \item[Why it matters:] Misdiagnosing the root cause leads to
    interventions (decoding fixes) that do not address visual failures.
  \item[Method:] Cosine similarity and Smooth Grad-CAM entropy to
    classify hallucinations as visual-origin vs.\ language-prior-origin.
  \item[Result:] 38--52\% of hallucinations are visual-origin; targeted
    visual fixes reduce POPE hallucination rate by 4.1 points and
    MSCOCO-Hallu rate by 37\% relative.
  \item[Evidence:] Three MLLM families; POPE, MME, MSCOCO-Hallu
    benchmarks; ablations of visual contrastive training and attention
    focusing.
\end{description}
\end{mdframed}

\section{The Language Prior Hypothesis and Its Limits}

The dominant account of object hallucination posits that MLLMs inherit
strong language priors from pre-training: the language model has seen
``refrigerator'' co-occurring with ``kitchen'' thousands of times and
will generate it even when the image shows an empty kitchen. Prior work
has proposed remedies at the language side: instruction tuning with
hallucination-negative examples, post-hoc decoding corrections, and
contrastive decoding between image-conditioned and image-unconditioned
outputs.

These language-side fixes improve hallucination metrics, but do not
fully close the gap with human performance. This paper asks whether
there is a second root cause that language-side fixes cannot address.

\section{Diagnosing Visual-Origin Hallucination}

The authors propose two diagnostic metrics that distinguish visual-origin
from language-prior-origin hallucination.

\textbf{Image-text cosine similarity} measures how well the visual
embedding aligns with the text of the generated output:
\[
  \text{sim}(v, t) = \frac{\mathbf{f}_v \cdot \mathbf{f}_t}
                          {\|\mathbf{f}_v\| \|\mathbf{f}_t\|}
\]
If the visual encoder correctly represents the image content,
$\text{sim}(v, t)$ should be high for faithful outputs. A low
$\text{sim}(v, t)$ for a hallucinatory output indicates that the
generated text does not correspond to the visual representation ---
suggesting the visual encoder failed.

\textbf{Smooth Grad-CAM entropy} measures the diffuseness of image
regions attended during generation:
\[
  H_{\text{Grad-CAM}} = -\sum_{r} p_r \log p_r
\]
where $p_r$ is the normalized Smooth Grad-CAM activation for image
region $r$. High entropy means the model fails to focus on
task-relevant objects. For faithful outputs, entropy is low; for
visual-origin hallucinations, entropy is systematically elevated.

Using these two metrics as a joint classifier, the authors assign each
hallucinated object instance to one of three categories:
\begin{itemize}[itemsep=1pt,parsep=0pt]
  \item \textbf{Language-prior-origin:} high sim, low entropy, object
    absent from image.
  \item \textbf{Visual-origin:} low sim or high entropy, regardless of
    whether the object could be predicted from language priors.
  \item \textbf{Ambiguous:} both factors elevated.
\end{itemize}

\section{Quantitative Decomposition}

Across three architectures (LLaVA-1.5, InstructBLIP, MiniGPT-4) on
MSCOCO-Hallu, the decomposition results are:

\begin{center}
\begin{tabular}{lccc}
\toprule
Model & Language & Visual & Ambig. \\
\midrule
LLaVA-1.5      & 41\% & 44\% & 15\% \\
InstructBLIP   & 38\% & 47\% & 15\% \\
MiniGPT-4      & 45\% & 40\% & 15\% \\
\midrule
Average        & 41\% & 44\% & 15\% \\
\bottomrule
\end{tabular}
\end{center}

The visual-origin fraction is consistently comparable to or larger than
the language-prior fraction, establishing it as a primary, not secondary,
cause.

\section{Targeted Remediation}

The authors propose two complementary fixes targeting the visual encoder.

\textbf{Visual contrastive training} pushes image embeddings toward
representations of faithful captions and away from unfaithful ones:
\[
  \mathcal{L}_{\text{contrast}} = -\log
  \frac{\exp(\text{sim}(v, t^+)/\tau)}
       {\exp(\text{sim}(v, t^+)/\tau) + \sum_j \exp(\text{sim}(v, t^-_j)/\tau)}
\]
where $t^+$ is the ground-truth caption and $t^-_j$ are negative captions
(hallucination-containing descriptions generated by the model).

\textbf{Attention focusing regularization} reduces visual entropy during
training by penalizing diffuse visual attention:
\[
  \mathcal{L}_{\text{focus}} = \lambda_2 \cdot H_{\text{Grad-CAM}}
\]
The full training objective combines both:
\[
  \mathcal{L} = \mathcal{L}_{\text{LM}}
              + \lambda_1 \mathcal{L}_{\text{contrast}}
              + \lambda_2 \mathcal{L}_{\text{focus}}
\]

\section{Experimental Results}

\begin{center}
\begin{tabular}{lcc}
\toprule
Model & POPE Acc. & Hallu.\ Rate \\
\midrule
LLaVA-1.5 (baseline)         & 85.3\% & 28.4\% \\
+ Language-side fix (ICD)     & 87.1\% & 24.9\% \\
+ Visual contrastive training & 88.6\% & 20.1\% \\
\textbf{+ Full method}        & \textbf{89.4\%} & \textbf{17.8\%} \\
\bottomrule
\end{tabular}
\end{center}

The full method reduces hallucination rate by 37.3\% relative on
MSCOCO-Hallu and improves POPE accuracy by 4.1 points, without
degrading object recognition or VQA performance.

\section{Ablations and Interpretation}

\textbf{Contrastive vs.\ focus ablation:} Contrastive training accounts
for 65\% of the hallucination reduction; attention focusing contributes
the remaining 35\%.

\textbf{Cross-architecture AUC:} The diagnostic metrics predict
visual-origin hallucinations with AUC 0.78--0.83 across all three
families, confirming the finding generalises beyond a single model.

\textbf{No regression on downstream tasks:} VQAv2 and ScienceQA scores
are unchanged, confirming the intervention is targeted and does not
erode general multimodal capability.

\textbf{Diagnostic threshold sensitivity:} The 41--44\% visual-origin
estimate shifts by at most $\pm$6 points across threshold choices,
showing the result is robust to the classification boundary.

\section*{Reference}

Peiyang Xu, Xiaopei Zhu, Jun Zhu, and Xiaolin Hu.
\textbf{Beyond Language Priors: Diagnosing and Fixing Visual-Origin
Hallucinations in Multimodal LLM.}
arXiv:2609.00231, August 2026.
\url{https://arxiv.org/abs/2609.00231}

\end{document}
